% Generated draft: requires mathematical and source review.
\documentclass[11pt]{article}
\usepackage[margin=25mm]{geometry}
\usepackage{fontspec}
\setmainfont{Latin Modern Roman}
\usepackage{amsmath,amssymb,mathtools,unicode-math}
\setmathfont{Latin Modern Math}
\usepackage{enumitem,needspace,xurl,hyperref}
\setlength{\parindent}{0pt}
\setlength{\parskip}{5pt}
\setlength{\emergencystretch}{3em}
\newcommand{\N}{\mathbb{N}}
\newcommand{\Z}{\mathbb{Z}}
\newcommand{\Q}{\mathbb{Q}}
\newcommand{\R}{\mathbb{R}}
\newcommand{\C}{\mathbb{C}}
\newcommand{\F}{\mathbb{F}}
\newcommand{\A}{\mathbb{A}}
\newcommand{\PP}{\mathbb{P}}
\newcommand{\E}{\mathbb{E}}
\newcommand{\eps}{\varepsilon}
\newcommand{\dd}{\,\mathrm{d}}
\begin{document}
\section*{UnsolvedMath ID: 700804}
\textbf{Source problem code:} catalogue entry
\subsection{Definitions and mathematical statement}
Let \(p\) be a prime. Write \(\mathcal K_p\) for the Kostrikin variety described in the source as the variety of locally finite groups of exponent \(p\). Here locally finite means that every finitely generated subgroup is finite, and the exponent condition means that \(g^p=1\) for every group element \(g\). Let \(K(p)\) be free in \(\mathcal K_p\) on a countably infinite set of generators: every assignment of these generators to elements of a group in \(\mathcal K_p\) extends uniquely to a homomorphism.

Under the usual lower central series interpretation of the source's phrase ``associated Lie ring,'' put
\[
\gamma_1(K(p))=K(p),\qquad \gamma_{n+1}(K(p))=[\gamma_n(K(p)),K(p)],
\]
\[
L=L(K(p))=\bigoplus_{n\geq 1}\gamma_n(K(p))/\gamma_{n+1}(K(p)).
\]
The addition on each summand is induced by its abelian quotient group structure, and the Lie bracket is induced by group commutators:
\[
[x\gamma_{i+1}(K(p)),y\gamma_{j+1}(K(p))]
=[x,y]\gamma_{i+j+1}(K(p))
\quad(x\in\gamma_i(K(p)),\ y\in\gamma_j(K(p))).
\]
The source does not explicitly specify this filtration; the formulation below uses this conventional interpretation and requires its confirmation against the original source.

Use ordinary Lie ring identities, without a grading restriction on substitutions. More precisely, let \(\mathfrak F\) be the free Lie ring over \(\Z\) on variables \(x_1,x_2,\ldots\), and let
\[
\operatorname{Id}(L)=\{f\in\mathfrak F: f(a_1,\ldots,a_m)=0\text{ for every substitution of its variables by elements of }L\}.
\]
For \(S\subseteq\mathfrak F\), write \(\langle S\rangle_T\) for the smallest Lie ideal of \(\mathfrak F\) that contains \(S\) and is invariant under every endomorphism of \(\mathfrak F\). Thus an identity follows from \(S\) when its polynomial belongs to \(\langle S\rangle_T\). A multilinear polynomial is a Lie polynomial homogeneous of degree one in each of its participating variables. A Lie ring is relatively free if it is free on some set in a variety of Lie rings defined by identities; equivalently, it has the universal mapping property for arbitrary assignments of its free generators into members of that variety.

For each prime \(p\), the original questions are:
\begin{enumerate}
\item[(a)] Is \(L(K(p))\) a relatively free Lie ring?
\item[(b)] Do all identities of \(L(K(p))\) follow from its multilinear identities? In the notation above, is
\[
\operatorname{Id}(L)=\left\langle\{f\in\operatorname{Id}(L):f\text{ is multilinear}\}\right\rangle_T?
\]
\item[(c)] Do all identities of \(L(K(p))\) follow from finitely many of its identities? That is, do there exist \(r\geq 0\) and \(f_1,\ldots,f_r\in\operatorname{Id}(L)\) such that
\[
\operatorname{Id}(L)=\langle f_1,\ldots,f_r\rangle_T?
\]
\end{enumerate}
These are separate questions with \(p\) arbitrary; no uniform finite basis independent of \(p\) is requested in (c).

The supplied record retains (a) and (c) in the main section and places (b) in the Archive of Solved Problems. Its archive entry reports a negative answer to (b): for \(p=7\), two relations of total weight \(29\), in two generators and of respective multiweights \((14,15)\) and \((15,14)\), hold in \(L(K(7))\) but are not consequences of multilinear relations. Here the stated multiweights count the occurrences of the two generators, and their sum is the total weight. This is the supplied historical resolution, not an independently verified conclusion. The record does not supply the two relations or explain explicitly how its wording about relations in generators implements the assertion about identities under arbitrary substitutions; that connection requires source review.
\subsection{Short English statement}
For each prime \(p\), form the associated Lie ring of the countably generated free group in the Kostrikin variety of locally finite groups of exponent \(p\). Ask whether this Lie ring is relatively free, whether its identities follow from multilinear identities, and whether its identities have a finite basis. The supplied record labels the problem partially solved: it reports a negative answer to the multilinear question at \(p=7\), while retaining the other two questions. The filtration convention and the reported passage from relations to identities require human review.
\subsection{Sources}
\begin{itemize}
\item Local UnsolvedMath problem registry, permanent ID \(700804\), exact source code \texttt{catalogue entry}, title ``Kourovka Notebook Problem 11.112.'' Supplied upstream revision: \texttt{d371ee917ac9647fbb06cd6fc529deec6f3dda82}.
\item E. I. Khukhro, proposer, 1990. \emph{The Kourovka Notebook (Unsolved Problems in Group Theory)}, 11th issue (1990), Problem 11.112; supplied record states that it was updated in the 21st issue (2026), arXiv:1401.0300v48. The supplied collection description names E. I. Khukhro and V. D. Mazurov as editors.
\item Supplied Archive of Solved Problems attribution for part (b): E. O'Brien and M. R. Vaughan-Lee, \emph{Int. J. Algebra Comput.} \textbf{12} (2002), 575--592. No article title or URL is supplied.
\item Further supplied archive attribution: M. F. Newman and M. R. Vaughan-Lee, \emph{Electron. Res. Announc. Amer. Math. Soc.} \textbf{4}, no. 1 (1998), 1--3. No article title or URL is supplied.
\item The supplied status review is dated \texttt{2026-10-16}, labels the record \texttt{partially\_solved}, and gives its evidence as ``published.'' Its note assigns solved parts to the archive and remaining parts to the main section. This is supplied dated curation, not an independently verified mathematical or literature assessment. The record's creation and update timestamps are both \texttt{2026-10-16T14:00:00Z}.
\end{itemize}
\end{document}
